% Intermediate summary matrix for Q7. May later be condensed, moved to an appendix, or redrawn as a figure.
\begin{table}[tbp]
\centering
\footnotesize
\singlespacing
\caption{Working summary for Q7: temporal predictability. Hazard and risk allocation are theory; lab timing effects use artificial sounds; field work addresses within-bout structure, not interval distributions.}
\label{tab:summary-q7}
\begin{tabular}{@{}p{1.35in}p{2.55in}p{1.2in}p{0.9in}@{}}
\toprule
\textbf{Factor} & \textbf{Finding} & \textbf{Stimulus class} & \textbf{Status} \\
\midrule
Risk allocation & Less vigilance per alarm when high-risk periods are frequent & Theory & Untested acoustically \\
\addlinespace[2pt]
Regular vs irregular ISI & Regular timing can promote habituation more than irregular at matched mean & Startle / tones & Strong (artificial) \\
\addlinespace[2pt]
Mean interval length & Short- and long-term habituation respond differently & Startle tones & Supporting \\
\addlinespace[2pt]
Hazard / distribution shape & Constant-hazard (exponential) predicted to preserve responding most & Hypothesis & Open \\
\addlinespace[2pt]
Prediction (DA recovery) & Unexpected noise slows NAc dopamine recovery & Aversive noise & Mechanistic bridge \\
\addlinespace[2pt]
Within-bout rate / structure & Call rate and repetition alter vigilance; age-dependent nulls exist & Natural sciurid calls & Strong (bout-level) \\
\addlinespace[2pt]
Natural calls $\times$ schedule & Interval distributions compared with natural calls & --- & Gap \\
\bottomrule
\end{tabular}
\end{table}
